\documentclass[10pt]{article}
\usepackage[margin=20mm]{geometry}
\usepackage{amsmath,amssymb,mathtools,fontspec,unicode-math,xurl,hyperref}
\IfFontExistsTF{DejaVu Serif}{\setmainfont{DejaVu Serif}}{\setmainfont{Latin Modern Roman}}
\setmathfont{latinmodern-math.otf}
\setlength{\parindent}{0pt}\setlength{\parskip}{6pt}\sloppy
\title{A Technical Review of
On the tightness of the maximum of branching Brownian motion in random environment}
\author{Initial scientific review: M25-0494}
\date{29 October 2026}
\begin{document}
\maketitle
Abstract: The acquired journal version was read, including its proof of quenched tightness. Several intermediate displays need correction, and the uniform time in Corollary 3.6 is not justified by its written choice of parameters. Explicit replacements are supplied under the existing assumptions. The overall assessment is Major: the written argument does not establish a uniform temporal shift essential to the central tightness proof. A separate comparison argument repairs this core gap under the existing assumptions; no false final tightness theorem is demonstrated. The remaining display and normalization defects are Minor. The review distinguishes the quenched conclusion, the annealed central limit discussion and the numerical nonconvergence conjecture.\par
\section*{Identity, assumptions and scope}
The reviewed paper is by Jiří Černý, Alexander Drewitz and Pascal Oswald, Annals of Probability 53 (2025), 509–543, DOI `10.1214/24-AOP1713`. The frozen, university-hosted journal PDF has 35 pages and SHA256\par
0d0bb57b14898a454c09a7b4ffb1d781fb2861ed05ffdfc3041b41e629f5e8b7\par
Below, printed journal page numbers are used; subtract 508 for PDF page numbers. The entire body, pp. 509–540, was read, as were the bibliographic entries on pp. 541–543. The corresponding arXiv record is 2212.12390v2, dated 24 June 2024; it is not substituted for the acquired journal file.\par
Theorem 2.1, p. 513, assumes spatial stationarity, the summable \(\psi\)-mixing condition (2.2), almost-sure local Hölder continuity, and \(0<\mathrm{ei}<\mathrm{es}<\infty\). Offspring numbers are at least one, their mean \(\mu\) exceeds one, and their second moment is finite. The additional ballistic condition \(v_0>v_c\) is essential for the compact negative tilting-parameter interval. The conclusion is tightness for almost every fixed environment under the quenched law; it does not give an environment-uniform radius or convergence in distribution. Figure 1 is expressly a discrete-space numerical motivation for a nonconvergence conjecture, not its proof. No author simulation was executed.\par
\section*{Independent checks and proof dependencies}
The normalization in Remark 3.2 is legitimate even when \(1<\mu<2\): \(\mu+p_1-2=\sum_{k\ge2}(k-2)p_k\ge0\). The new probabilities sum to one, have mean two, and preserve the rates of every genuine branching event. Equivalently, the transformed reaction multiplied by \((\mu-1)\xi\) equals the original reaction multiplied by \(\xi\). Exact symbolic checks used \((p_1,p_2,p_3)=(2/5,3/10,3/10)\), giving new probabilities \(1/3\) each. The Taylor lower bound (6.7) was also checked algebraically.\par
The tilted diffusion calculation on pp. 522–525 has \(b=(\log Z)'\), \(b'=-2(\zeta+\eta)-b^2\) and \[\sqrt{2|\eta|}\le b(x)\le\sqrt{2(\mathrm{es}-\mathrm{ei}+|\eta|)}.\] This gives a bounded Lipschitz drift and the stopped density in (4.22). The drift comparison can be proved directly on a common Brownian path: \(x+B_t+\underline v t\le X_t\le x+B_t+\overline v t\). The constant-killing diagnostic \(Z(x)=e^{\sqrt{2a}x}\) satisfies \(Z''/2-aZ=0\); it checks signs, not the nondegenerate-environment theorem.\par
The mixing assumption enters the imported concentration estimates in Section 5 even though it is not used directly in the local PDE steps. Endpoint uniformity requires the union bounds and compact speed set in Lemma 5.3; the compact set must lie strictly above \(v_c\). For Lemma 6.1, perturbations need first be applied at one fixed sufficiently large displacement; monotonicity in the target then extends the result to every larger displacement. Applying a sublinear perturbation estimate directly to arbitrary unbounded displacements would be unjustified.\par
A useful independent check of the constants in Lemma 6.4 keeps \(s=\sqrt{\operatorname{Var}_{\widetilde\eta}H_y}\) separate from \(d=\eta-\widetilde\eta>0\), so that \(\widetilde\sigma=ds\) and the integration cutoff is \(\ell=L/s\). From the stated uniform variance and interval normal-approximation estimates, \(s\le b\sqrt{vt}\) and error at most \(C_{17}/\sqrt{vt}\). Choose \(L\ge20bC_{17}\). For fixed \(L\), sufficiently large \(t\) gives the numerator and denominator bounds \[N\le\frac{2\ell e^{-dL}}{1-e^{-dL}},\qquad
D\ge\frac{\ell}{20}e^{-dL/2}.\] Thus \(N/D\le40e^{-dL/2}/(1-e^{-dL})\). With \(dL\ge12\), the computed bound is \(0.099151<1/4\). This tracks the dependence as \(d\) becomes small, instead of treating the lower bound in (6.34) as uniform in \(L\). It verifies the final comparison conditional on the stated variance and normal-approximation inputs; their external proofs were not reproduced.\par
Finally, the Sturm comparison on p. 540 transfers a nonnegative difference at a negative spatial point to the origin. The quantile law of large numbers places that point in the range of Lemma 6.1. Bounded time intervals are harmless: bounded branching rates and finite offspring mean give finitely many particles on compact time intervals almost surely, with continuous trajectories and maxima. This supplies the compact-time part omitted from the asymptotic argument.\par
\section*{Problem 1: factors and signs in Proposition 5.1}
Severity: Minor; local repairs preserve both perturbation conclusions.\par
\subsection*{The missing potential-shift factor}
Equation reference: eq:shift. Equation (5.24), p. 531, applies Lemma 5.6 to two different times but omits \(e^{\mathrm{es}h}\). Define \[U_\xi(t;x,y)=E_x[e^{\int_0^t\xi(X_s)\,ds};X_t\ge y].\] Since \(\zeta=\xi-\mathrm{es}\), the exact identity is \(U_\xi(t;x,y)=e^{\mathrm{es}t}U_\zeta(t;x,y)\). For \(v=(y-x)/t\), \(v'=(y-x)/(t+h)\), (5.23) therefore gives \[
\frac{U_\xi(t+h;x,y)}{U_\xi(t;x,y)}
\le e^{\mathrm{es}h}C_{12}^{\,2}
       \frac{Y^\approx_{v'}(x,y)}{Y^\approx_v(x,y)} .\] The omitted factor cannot be absorbed into a fixed prefactor in (5.24). Indeed the opposite halves of (5.23) give a lower bound with \(e^{\mathrm{es}h}C_{12}^{-2}\) times the same positive ratio. Fix \(v\) in the interior of \(V\), set \(x=0,y=vt\), take \(A>v\), and use \(h=\sqrt t\), \(\delta=1/2\). For large \(t\), all the proposition’s endpoint and speed conditions hold. The printed bound would force \(e^{\mathrm{es}\sqrt t}\le C_{12}^{4}\), a contradiction.\par
The counterevidence against a false proposition is decisive: its claimed conclusion already permits \(C_1e^{C_1h}\). Restoring [eq:shift] (\#eq:shift) and increasing \(C_1\) absorbs the factor without changing an assumption or the conclusion.\par
\subsection*{The negative spatial exponent}
Equations (5.30)–(5.31), pp. 531–532, retain only an absolute \(O(h)\) estimate, which alone cannot yield the asserted \(e^{-h/C_2}\) decay. With the same tilting parameter \(\eta_* \in\Delta\Subset(-\infty,0)\) in the first difference of (5.30), its actual signed value is \[\log Z^{\zeta,\eta_*}_{y,y+h}
 \le-\sqrt{2|\eta_*|}\,h\le-a h\] for \(a>0\) uniform on \(\Delta\). This follows from (4.17)–(4.18) and \(\zeta\le0\). Keeping the normal-approximation prefactor also contributes an additive bounded constant in the logarithm. The corrected logarithmic bound is \[O(1)-a h+O(h^2/t).\] Because \(h\le t\delta(t)\) with \(\delta(t)\to0\), it yields the required negative exponent for large \(t\). For \(0\le h<1\), monotonicity of the target event and an enlarged prefactor suffice. These corrections repair the displayed derivation feeding Lemma 6.1; they do not demonstrate a false final theorem.\par
\section*{Problem 2: continuation after a tilted hitting time}
Severity: Minor. In (6.17), p. 535, the change-of-measure density stops at \(H_y\), but the next equality uses \(p_y^{\zeta,\eta}(t-H_y)\), the endpoint probability under the tilted continuation. This continuation was not changed by a density measurable at \(H_y\).\par
On the event \(A=\{H_y\in[t-K,t],\,T_{y,t}\ge t-K\}\), the barrier condition is determined by the path up to \(H_y\): the first barrier crossing occurs no later than \(H_y\). Conditional on this path, the original Brownian continuation satisfies \[P_y(X_{t-H_y}\ge y)=\tfrac12\] for \(H_y<t\); the endpoint \(H_y=t\) is a null event and would give one. The correct stopped-density identity uses this Brownian probability, so it directly yields the final factor \(1/2\) required in (6.17).\par
This is a false intermediate equality, rather than evidence against the lower bound. Under the tilted continuation the positive drift gives \(p_y^{\zeta,\eta}(s)\ge\Phi(\sqrt{2|\eta|s})>1/2\) for \(s>0\). For the constant-drift diagnostic with drift two at time one, the computed probability is \(0.977249868\), compared with \(1/2\) for Brownian motion. The repair is to condition before tilting and retain the Brownian continuation. The subsequent proof of Lemma 6.2 uses only the unchanged \(1/2\) bound.\par
\section*{Problem 3: uniformity of the growth time}
Severity: Major; a core proof gap with a verified repair. Corollary 3.6, pp. 520–521, requires one time \(T\) for all sufficiently large \(y\) and any \(t\) with \(w^y(t,0)=\varepsilon\). Its proof puts \(t=\tau_y^\varepsilon+s_0\), chooses \(s_1\) using \(s_0\), and sets \(T=s_1-s_0\). Both \(s_0\) and the compact interval on which convergence is used may vary with \(y,t\). Local uniform convergence on fixed compact intervals does not justify this uniform choice. Restricting to first hitting times would set \(s_0=0\), but the stated corollary and its application allow any such \(t\). Theorem 2.1 uses this uniform shift in (7.4) before the Sturm comparison, so the gap affects a necessary step of the central tightness derivation. The Major grade concerns this invalid passage of quantifiers; the availability of the repair below does not validate the written argument or imply that the final theorem is false.\par
Here is a direct repair under the existing assumptions, in the normalization \(\mu=2\). Write \(q=1-p_1>0\). Formula (3.1) implies \(F(w)\ge q w(1-w)\) on \([0,1]\). By the heat-kernel Duhamel formula, for \(t\ge1\), \[\|\partial_xw^y(t,\cdot)\|_\infty
\le C:=\sqrt{2/\pi}\,(1+2\mathrm{es}),\] using \(0\le w,F\le1\). At a time with \(w^y(t,0)=\varepsilon\), therefore, \(w^y(t,x)\ge\phi(x):=(\varepsilon-C|x|)_+\). Reaction is nonnegative, so one time unit later the profile dominates \(P_1\phi\), strictly positive on every compact interval.\par
Set \(c=\mathrm{ei}\,q>0\), choose \(R\) so that \(\kappa:=\pi^2/(8R^2)<c\varepsilon/2\), and let \(\psi(x)=\cos(\pi x/(2R))\) on \([-R,R]\). Take \(0<a_0<\min\{1-\kappa/c,\min_{[-R,R]}P_1\phi\}\) and solve \[a'=a(c-\kappa-ca),\qquad a(0)=a_0.\] Then \(a(s)\psi(x)\) is a subsolution on this interval with zero boundary values. Indeed, the reaction-diffusion right side minus its time derivative, divided by \(\psi\), is \(ca(s)^2(1-\psi)\ge0\). Comparison applies because \(\xi F(w)\ge c w(1-w)\). Moreover \(a(s)\uparrow1-\kappa/c>1-\varepsilon/2\). One deterministic time, depending only on \(\varepsilon,\mathrm{ei},\mathrm{es},p\), makes the value at the origin at least \(1-\varepsilon\) for all subsequent times. The large-\(y\) regime ensures the original \(t\ge1\). This proves the corollary’s uniform claim and its use in (7.4), without strengthening the environment assumptions. The symbolic subsolution residual was checked independently.\par
\section*{Problem 4: introductory normalization displays}
Severity: Minor. Equation (1.5), p. 510, subtracts the survival probabilities in the wrong order. Since \(x\mapsto w(t,x)=P(M(t)\ge x)\) is nonincreasing, \(w(t,m+r)-w(t,m-r)\le0\), whereas \(1-\varepsilon>0\). Reverse the subtraction to express the probability of the central interval.\par
Remark 2.3, pp. 513–514, first states convergence of \((M(nt)-v_0nt)/\sqrt{\sigma^2n}\) to standard Brownian motion, but later assigns the same standardized marginal an \(N(0,\sigma^2)\) limit. At fixed \(t\) the standard-Brownian marginal is \(N(0,t)\), and at \(t=1\) it is \(N(0,1)\). Scaling by \(\sqrt n\) instead yields variance \(\sigma^2t\). These are direct internal normalization inconsistencies. The main quenched tightness proof uses neither erroneous display.\par
\section*{Assessment and limits}
The corrected arguments preserve the theorem’s range and its dependence on \(v_0>v_c\), summable mixing, local regularity, positive bounded branching rates and finite offspring second moment. The zero-number theorem, quenched law of large numbers, mixing concentration and uniform normal-approximation inputs remain external dependencies; their full original PDFs were not inspected. Closed versus open Heaviside values at one point do not affect positive-time Brownian kernels, though that convention should be stated when applying the lower-semicontinuous-data formulation of Proposition 3.1.\par
The author publication page and arXiv metadata were checked. No separately verified correction addressing these specific displays was found in the bounded search; this is not a claim of exhaustive absence. Source, independent calculations, compiler and retained dependency payloads, and actual page inspections are bound in the support inventory. Scientific decision remains pending and human reviews are empty.\par
9 J. Černý, A. Drewitz and P. Oswald, On the tightness of the maximum of branching Brownian motion in random environment, Ann. Probab. 53 (2025), 509–543. DOI `10.1214/24-AOP1713`. Acquired journal PDF. arXiv, record for 2212.12390, https://arxiv.org/abs/2212.12390. J. Černý, publication list, University of Basel, https://probability.dmi.unibas.ch/cerny/publications.html.\par
\end{document}
